\Introduction

Актуальность темы. В современной практике обеспечения качества сложных программных комплексов проведение нагрузочного тестирования требует создания стабильной и управляемой среды~\cite{gost-56920-2024, kulikov-testing}. Одной из ключевых проблем при организации такой среды является управление инфраструктурой имитационных сервисов (заглушек) -- специализированного программного обеспечения, воспроизводящего поведение реальных внешних компонентов системы. Применение имитационных сервисов позволяет изолировать тестируемую систему от реальных внешних зависимостей, обеспечивая воспроизводимость тестовых сценариев и возможность моделирования различных режимов работы, в том числе нештатных. В условиях распределённых архитектур количество задействованных заглушек существенно возрастает, а их ручное сопровождение становится трудоёмким, что обусловливает необходимость автоматизированных средств управления.

В условиях реализации политики импортозамещения в Российской Федерации актуальность приобретает разработка отечественных инструментов управления инфраструктурой, способных заменить зарубежные программные решения и обеспечить совместимость с сертифицированными отечественными операционными системами~\cite{burenin-astra}.

Цель работы -- проектирование и разработка автоматизированной информационной системы для централизованного управления имитационными сервисами в среде нагрузочного тестирования и мониторинга их состояния.

Задачи исследования:
\begin{compactlist}
  \item Провести анализ предметной области и существующих подходов к управлению имитационными средами в процессах обеспечения качества ПО.
  \item Спроектировать архитектуру программного комплекса, обеспечивающую гибкое масштабирование и работу в различных режимах.
  \item Обосновать выбор технологического стека и программных средств реализации системы.
  \item Разработать алгоритмы динамического проксирования трафика, механизмы контроля интенсивности запросов и программные интерфейсы для автоматизации управления.
  \item Реализовать систему хранения конфигураций и контроля состояний исполняемых файлов.
  \item Провести тестирование разработанного комплекса в среде ОС Astra Linux и оценить эффективность реализованных механизмов мониторинга и управления.
\end{compactlist}

Объект исследования -- инфраструктура запуска и эксплуатации имитационных сервисов (заглушек) в среде нагрузочного тестирования.

Предмет исследования -- механизмы и алгоритмы автоматизированного управления и контроля состояния имитационных сервисов в среде нагрузочного тестирования.

Научная новизна работы заключается в разработке алгоритма централизованного управления распределённой средой, поддерживающего индивидуальную настройку ограничений интенсивности запросов (Rate~Limiting) для каждого имитационного сервиса в отдельности, а также в реализации агентно-независимой схемы взаимодействия с удалёнными узлами посредством протоколов SSH/SCP (Secure Shell, Secure Copy Protocol) без установки дополнительного программного обеспечения на целевые серверы.

Практическая значимость работы состоит в создании программного инструмента, позволяющего отказаться от использования зарубежных систем управления, упростить эксплуатацию инфраструктуры имитационных сервисов и обеспечить её интеграцию со стандартными средствами мониторинга (Prometheus).

Положения, выносимые на защиту:
\begin{compactlist}
  \item Архитектура программного комплекса, поддерживающая гибкое масштабирование системы от локального запуска на одном хосте до формирования распределённого кластера.
  \item Механизм динамического проксирования запросов, реализующий прозрачную маршрутизацию входящего трафика к целевым имитационным сервисам.
  \item Алгоритм контроля интенсивности запросов (Rate~Limiting), обеспечивающий стабильность работы имитационной среды при пиковых нагрузках.
  \item Методика формирования и экспорта метрик производительности для непрерывного мониторинга состояния инфраструктуры заглушек.
\end{compactlist}

Структура работы. Выпускная квалификационная работа включает введение, три главы основного текста, заключение, список использованных источников и приложения~\cite{gost-7-32-2017}.

Во введении обоснована актуальность исследования, поставлены цель и задачи, определены объект и предмет работы, указаны научная новизна и практическая значимость полученных результатов.

В первой главе проводится анализ предметной области, исследуются существующие аналоги и методы управления имитационными средами. На основе сравнительного анализа обосновывается выбор технологий и формируются требования к системе.

Во второй главе описывается проектирование архитектуры системы, структуры хранения данных и логики работы основных модулей. Приводятся схемы взаимодействия компонентов и алгоритмы реализации функций проксирования и управления кластером.

В третьей главе рассматриваются вопросы практической реализации, настройки и эксплуатации системы в среде ОС Astra Linux. Приводятся результаты тестирования, подтверждающие корректность работы разработанных механизмов управления и мониторинга.

В заключении сформулированы основные выводы, подтверждающие решение поставленных задач и достижение цели работы.